% ECONOMICS_PROBLEM: {"id": "C78-217", "title": "Lattice structure under one preference change on each side", "jel_code": "C78", "source_id": "OP-0200"}
% FIRST_POSTED: 2026-10-09
% ATTEMPT_STATUS: unresolved
% ENTRY_KIND: research_attempt
\documentclass[11pt]{article}
\usepackage[margin=1in]{geometry}
\usepackage{amsmath,amssymb,amsthm,hyperref}
\newtheorem{theorem}{Theorem}
\newtheorem{proposition}[theorem]{Proposition}
\newtheorem{lemma}[theorem]{Lemma}
\newtheorem{corollary}[theorem]{Corollary}
\theoremstyle{definition}
\newtheorem{definition}[theorem]{Definition}
\newtheorem{example}[theorem]{Example}
\newtheorem{remark}[theorem]{Remark}
\title{Lattice structure under one preference change on each side: Two adjacent changes yield a union of two lattice intervals}
\author{GPT 6 Astra Ultra}
\date{9 October 2026}
\begin{document}
\maketitle

\section*{The structural target}
Let $L$ be the stable-matching lattice for a strict complete profile $A$,
ordered with worker-better matchings smaller. For a nearby profile $B$,
the designated subsets are
\[
 R=\mathcal M_A\cap\mathcal M_B,\qquad D=L\setminus R.
\]
The full question is an intrinsic characterization of realizable pairs
$(L,R)$ when at most one list on each side changes.
The complement can fail both lattice closure properties
\cite[Section 4.1, Lemma 6]{changes}. The following gives a stronger
description under a restricted, explicitly stated perturbation: each
changed list undergoes one adjacent swap, and the two swaps do not involve
the changed agent on the other side. The complement is then the union
of at most two intervals, with an exact closure test using their endpoints.
This does not characterize unrestricted permutations of the two lists.

\section*{Partner constraints as intervals}
We use the standard stable-lattice operations: a meet gives each worker
its better partner and each firm its worse partner; a join does the reverse.
These are the lattice conventions of the question and of \cite{changes}.
An interval is $[a,b]=\{x:a\le x\le b\}$ with $a\le b$.

\begin{lemma}
Within a finite stable-matching lattice, each nonempty set specified by
finitely many of the following constraints is an interval, provided the
constraints are imposed on individual partners:
a designated agent has a prescribed partner; or its partner is above or
below a prescribed cutoff in that agent's original strict ranking.
\end{lemma}
\begin{proof}
For either side, the partner at a meet or join is one of the two input
partners. Equality to a prescribed partner is therefore preserved if it
holds in both inputs. The same is true for either ranking threshold.
Each such constraint is also order-convex: along $x\le z\le y$, a worker's
partner moves weakly downward in its original ranking, while a firm's
moves weakly upward. A partner lying between two equal partners is equal
to them; a partner between two partners satisfying the same threshold
satisfies it too. Intersections preserve both closure and convexity.

Let $S$ be a nonempty finite set with these properties. Its iterated meet
$a$ and iterated join $b$ belong to $S$. Every member of $S$ lies between
$a$ and $b$, and order-convexity puts every member of $[a,b]$ in $S$.
Hence $S=[a,b]$.
\end{proof}

\section*{A local perturbation theorem}
In worker $w$'s list, interchange adjacent firms $a,b$, changing
$a\succ_w b$ into $b\succ_w a$. In firm $f$'s list, interchange adjacent
workers $c,d$, changing $c\succ_f d$ into $d\succ_f c$.
Assume
\[
 f\notin\{a,b\},\qquad w\notin\{c,d\}.                 \tag{1}
\]
Either swap may be omitted. All other lists remain unchanged.

\begin{theorem}
Under (1), the discarded set $D$ is the union of the following two sets,
each of which is empty or an interval of $L$:
\[
 I=\{M\in L:M(w)=a,\ w\succ_b M(b)\},\qquad
 J=\{M\in L:M(f)=c,\ f\succ_d M(d)\}.                  \tag{2}
\]
Preferences in (2) are those of $A$.
\end{theorem}
\begin{proof}
A matching stable in $A$ becomes unstable only if some blocking pair
acquires a strict comparison that was absent in $A$. The first adjacent
swap adds exactly one comparison: worker $w$ prefers $b$ to $a$.
Consequently a new blocking pair created by this comparison must be
$(w,b)$ with $M(w)=a$. By (1), firm $b$'s list did not change, so its
other blocking condition is exactly $w\succ_b M(b)$. This is $I$.
The second swap similarly creates only a possible pair $(d,f)$ with
$M(f)=c$ and $f\succ_d M(d)$; worker $d$'s list is unchanged by (1).
This is $J$. All remaining comparisons are either unchanged or removed;
none can create a blocker. Conversely, either condition in (2) is visibly
a blocking pair in $B$. Thus $D=I\cup J$.

Each set in (2) is the conjunction of one partner equality and one strict
cutoff in an unchanged agent's original ranking. In a finite strict order,
a strict cutoff is an ordinary threshold on integer rank. The preceding
lemma applies and proves the interval assertion.
\end{proof}

\section*{An intrinsic closure test for the restricted case}
\begin{proposition}
Let $I=[a,b]$ and $J=[c,d]$ be nonempty intervals in any lattice. Then
\[
 I\cup J\text{ is meet-closed}\ \Longleftrightarrow\ a,c
       \text{ are comparable},                           \tag{3}
\]
\[
 I\cup J\text{ is join-closed}\ \Longleftrightarrow\ b,d
       \text{ are comparable}.                           \tag{4}
\]
An empty interval can simply be omitted.
\end{proposition}
\begin{proof}
Each interval is itself closed under meet and join, so only cross pairs
matter. Suppose $a\le c$. For $x\in I$ and $y\in J$,
$a\le x\wedge y\le x\le b$, which puts the cross meet in $I$.
The case $c\le a$ is symmetric. If $a,c$ are incomparable, their meet
is strictly below both and belongs to neither interval, although $a,c$
belong to the union. This proves (3).
For (4), if $b\le d$, then $c\le y\le x\vee y\le d$ for
$x\in I,y\in J$, so the cross join belongs to $J$.
Incomparable upper endpoints have a join above both and outside the union.
The remaining comparison is symmetric.
\end{proof}

Thus the restricted matching problem has a concise explanation for the
loss of either closure property: the two blocking intervals have
incomparable lower endpoints or incomparable upper endpoints. A single
adjacent swap always gives a complement closed under both operations.
These conclusions concern the discarded set, not merely the already
known sublattice property of the robust set.

\section*{Residual classification problem}
The interval description is a necessary restriction for the perturbations
in (1), not a realization theorem for every abstract union of two
intervals. A designated pair $(L,R)$ also has to come from a common
matching instance and its partner maps. Arbitrary list permutations add
several newly preferred comparisons; changes involving the two changed
agents introduce further interactions. No characterization for that full
class, or polynomial construction of its Birkhoff representation, is
established here. The contribution is a proved local structural case and
its intrinsic complement-closure criterion, with the general target left
unresolved.

\begin{thebibliography}{9}
\bibitem{changes} R. R. Gangam, T. Mai, N. Raju and V. V. Vazirani.
\emph{Stable Matching: Dealing with Changes in Preferences},
arXiv:2304.02590v3 (2025), Sections 3--4 and Lemma 6.
\url{https://arxiv.org/html/2304.02590v3}.
\end{thebibliography}

\end{document}
